\documentclass[a4paper,landscape, 12pt]{extarticle}
\usepackage[left=0.5in, right=0.5in, top=0.4in, bottom=0.55in]{geometry}
\usepackage{graphicx} % Required for inserting images
\usepackage[normalem]{ulem}
\usepackage{tabto}
\usepackage{minted}
\renewcommand\ULthickness{1.0pt}   %%---> For changing thickness of underline
\setlength\ULdepth{1.3ex}%\maxdimen ---> For changing depth of underline
\usepackage{tikz}
\usetikzlibrary{graphs}
\usepackage{multicol}
\usepackage{array}
\usepackage{caption}
\usepackage{amsmath} 
\usepackage{fancyhdr}
\usepackage{lastpage}

\author{Juan R Ortiz}

% Page style with slide numbers
\pagestyle{fancy}
\fancyhf{}
\renewcommand{\headrulewidth}{0pt}
\fancyfoot[R]{\Large\thepage\ / \pageref{LastPage}}

% Define header command for consistent formatting
\newcommand{\slideheader}[2]{%
\noindent

\includegraphics[height=0.5in]{cpp_logo.png}%
\hspace{0.3in}%
{\Large \uline{#1 \hfill #2}}%
\par\vspace{0.4in}
}
% Standard content spacing - 0.6in base indent for all content
\newcommand{\contentstart}{%
\leftskip 0.6in \rightskip 0.6in
}
\newcommand{\contentend}{%
\leftskip 0in \rightskip 0in
}

% Standard font size for body content
\newcommand{\bodytext}{\LARGE}

\begin{document}
\begin{huge} 
%%%%%%%%%%%%%%%%%%%%%%%%%%%%%%%%%%%%%%%%
% SLIDE 1 - Title Slide
\slideheader{ECE5110 Numerical Modeling}{Unit 02}

\vspace{1.2in}
\begin{center}
{\Huge \textbf{Juan Ortiz}}

\vspace{0.4in}
{\LARGE Polynomial Interpolation for Modeling DC Motor Calibration}
\end{center}
%%%%%%%%%%%%%%%%%%%%%%%%%%%%%%%%%%%%%%%%
\newpage %Page 2
\slideheader{Unit 02: Polynomial Interpolation}{ECE5110 - Numerical Modeling}

\leftskip 0.7in \rightskip 0.7in 
\vspace{0.4in}
\noindent A manufacturing robot uses a small brushed DC motor to control conveyor belt speed. Before tuning the PID controller, a static model of the motor’s steady‑state response is needed. 

\vspace{0.2in}
This model must show how the motor reacts to different control voltages.

\vspace{0.4in}
\begin{center}
\includegraphics[height=2.5in]{unit02/manufacturing.png}%
\end{center}

\leftskip 0in \rightskip 0in 

%%%%%%%%%%%%%%%%%%%%%%%%%%%%%%%%%%%%%%%%

\newpage %Page 3
\slideheader{Unit 02: Polynomial Interpolation}{ECE5110 - Numerical Modeling}

\leftskip 0.7in \rightskip 0.7in 
\vspace{0.4in}
During testing, data points were collected at several fixed voltages between 0.5V and 5.0V. At each setting, the system was allowed to reach its steady‑state speed.

\vspace{0.6in}
\begin{table}[h!]
\centering
\begin{tabular}{|c|c|} % Defines two centered columns with vertical lines
\hline
\textbf{Voltage (V)} & \textbf{Speed (RPM)} \\ \hline
0.5 & 0 \\ \hline
1.0 & 80 \\ \hline
1.5 & 200 \\ \hline
2.0 & 380 \\ \hline
2.5 & 650 \\ \hline
3.0 & 1000 \\ \hline
3.5 & 1350 \\ \hline
4.0 & 1600 \\ \hline
4.5 & 1750 \\ \hline
5.0 & 1850 \\ \hline
\end{tabular}
\caption{Motor Steady-State Response Across Fixed Voltage Inputs}
\end{table}

\leftskip 0in \rightskip 0in 

%%%%%%%%%%%%%%%%%%%%%%%%%%%%%%%%%%%%%%%%
\newpage %Page 4
\slideheader{Unit 02: Polynomial Interpolation}{ECE5110 - Numerical Modeling}

\leftskip 0.7in \rightskip 0.7in 
\vspace{0.4in}
Given the testing data, we can now approximate a S-Curve that will pass through each of those points.

\vspace{0.2in}
\noindent \textbf{Lagrange Interpolation Method:}
\vspace{0.2in}

\noindent For each point ($x_k$, $y_k$), a "mini-curve" basis polynomial $L_k(x)$ is created, where it equals 1 at $x_k$ and 0 at all other data points. 
\vspace{0.2in}

\noindent The final S-Curve will be the SUM of all the mini-curves scaled by $y_k$:
\[\hspace{0.6in}
P(x) = y_0L_0(x) + y_1L_1(x) + ... + y_{n-1}L_{n-1}(x)
\]
\[\hspace{0.6in}
P(x) = \sum_{k=0}y_k\frac{\Pi(X-x_i)}{\Pi(x_k-x_i)}
\]

\leftskip 0in \rightskip 0in 

%%%%%%%%%%%%%%%%%%%%%%%%%%%%%%%%%%%%%%%%
\newpage % Page 5
\slideheader{Unit 02: Polynomial Interpolation}{ECE5110 - Numerical Modeling}

\leftskip 0.7in \rightskip 0.7in 
\vspace{0.4in}
\noindent In Python, we can compute the denominator by:
\vspace{0.2in}

\begin{center}
\begin{minipage}{0.75\linewidth} % adjust width as you like
\begin{minted}[fontsize=\huge, frame=single]{python}
diff = x[k] - np.delete(x, k)
denom[k] = np.prod(diff)
\end{minted}
\end{minipage}
\end{center}

\noindent \hspace{4.5in} equivalently: ${\Pi_{i \neq k}(x_k-x_i)}$
\vspace{0.2in}

\noindent Followed by building the numerator:

\begin{center}
\begin{minipage}{0.75\linewidth}
\begin{minted}[fontsize=\huge, frame=single]{python}
num = np.ones_like(X)
for i in range(n):
    if i == k: continue
    num *= (X - x[i])
\end{minted}
\end{minipage}
\end{center}

\noindent \hspace{4.5in} equivalently: ${\Pi_{i \neq k}(X - x_i)}$

\leftskip 0in \rightskip 0in

%%%%%%%%%%%%%%%%%%%%%%%%%%%%%%%%%%%%%%%%

\newpage % Page 6
\slideheader{Unit 02: Polynomial Interpolation}{ECE5110 - Numerical Modeling}

\leftskip 0.7in \rightskip 0.7in 
\vspace{0.4in}
\noindent Lastly, the code will build the interpolation:
\vspace{0.2in}

\begin{center}
\begin{minipage}{0.75\linewidth} % adjust width as you like
\begin{minted}[fontsize=\huge, frame=single]{python}
out += y[k] * (num / denom[k])

\end{minted}
\end{minipage}
\end{center}

\noindent \hspace{4.5in} equivalently: $\sum_{k=0}y_k\frac{\Pi(X-x_i)}{\Pi(x_k-x_i)}$
\vspace{0.2in}

\noindent This ensures the output matches the exact data point and fills in the curve smoothly for every value in between.

\leftskip 0in \rightskip 0in

%%%%%%%%%%%%%%%%%%%%%%%%%%%%%%%%%%%%%%%%

\newpage % Page 7
\slideheader{Unit 02: Polynomial Interpolation}{ECE5110 - Numerical Modeling}

\leftskip 0.7in \rightskip 0.7in 
\vspace{-0.4in}
\begin{center}
\[
\begin{aligned}
x &= [0.5, 1.0, 1.5, 2.0, 2.5, 3.0, 3.5, 4.0, 4.5, 5.0] \\
y &= [0, \; 80, \; 200, \; 380, \; 650, \; 1000, \; 1350, \; 1600, \; 1750, \; 1825]
\end{aligned}
\]
\end{center}

\begin{center}
\includegraphics[height=4in]{unit02/SCurve_Lagrange.png}%
\end{center}


\leftskip 0in \rightskip 0in

%%%%%%%%%%%%%%%%%%%%%%%%%%%%%%%%%%%%%%%%

\newpage % Page 8
\slideheader{Unit 02: Polynomial Interpolation}{ECE5110 - Numerical Modeling}

\leftskip 0.7in \rightskip 0.7in 
\vspace{0.4in}
\begin{itemize}
    \item The plot represents the static plant model (input $\rightarrow$ steady-state output) of the brushed DC motor.
    \item A deadzone appears at low voltages due to mechanical friction.
    \item The motor speed increases rapidly through the mid-voltage region.
    \item At higher voltages, the motor approaches saturation, where additional voltage results in only small increases in RPM.
\end{itemize}

\vspace{0.2in}
Based on these findings, we can tune a PID (Proportional Derivative Integral) controller to the characteristics of the specific brushed DC motor.

\leftskip 0in \rightskip 0in

%%%%%%%%%%%%%%%%%%%%%%%%%%%%%%%%%%%%%%%%

\newpage %End Page
\slideheader{Unit 02: Polynomial Interpolation}{ECE5110 - Numerical Modeling}

\leftskip 0.7in \rightskip 0.7in 
\vspace{2.5in}
\begin{center}
{\Huge \textbf{QUESTIONS?}}
\end{center}
\leftskip 0in \rightskip 0in 


%%%%%%%%%%%%%%%%%%%%%%%%%%%%%%%%%%%%%%%%


\end{huge}
\end{document}